\section{犯罪形态 - 犯罪既遂}
\begin{enumerate}
  \item 结果要求。要有实害结果。
  \item \zd{产生危险是故意犯罪的成立条件，造成实害结果是故意犯罪的既遂条件}。
  \item 行为对象的转移。\zd{可替代性}。
        \begin{enumerate}
          \item 人身权益。专属性，账一个一个算。\al{强奸 A 转而强奸 B}。
          \item 财产法益。不具有人身专属性，具有可替代性。\al{买椟还珠 \yjt 盗窃既遂}。
        \end{enumerate}
  \item 属于\zd{实行阶段行为导致}的，不能是预备阶段产生的实害结果。
  \item vs 未遂、中止。
        \begin{enumerate}
          \item \zd{\blx{都是终局形态（互斥）}}。
          \item 考试套路。1 个犯罪，弄 2 个时间点。容易犯的错，先判断第 2 个时间点。
          \item 做题画时间轴，\zd{\blx{从第 1 个时间节点往后看}}有没有形成终局形态。
          \item 判断标准。\zd{客观上犯罪行为彻底结束，主观上犯意要彻底消除}（这事儿翻片）。
          \item \al{甲欲杀乙，砍了 10 刀，认为必死无疑，离开现场。2 小时后想到刀还在现场，发现乙没死，送去医院。中止}。\al{打扫“战场”、离去、次日}。
          \item \al{想杀丈夫，灌醉，关好门窗，打开煤气，离开现场，在院子里等，2 小时后返回现场，发现人没死，送去医院。未遂}。没考。
        \end{enumerate}
  \item \zd{重复侵害行为}。新增考点。\al{5 发子弹，打了第 1 发后放弃}。
        \begin{enumerate}
          \item 个别考察说。第 1 枪未遂，第 2 枪中止。
          \item \zd{整体考察说}。第 1 枪后放弃，犯意没消除，没有形成终局形态，属于自动放弃 \yjt 中止。
        \end{enumerate}
  \item 多个行为才能达目的。\zd{喜欢考行为结构复杂的案件}。\al{诈骗罪、拐卖、拐骗、强奸罪}。
  \item \zd{掌握常见各个罪名的成立和既遂，互斥}，不要被多余描述干扰。
        \begin{enumerate}
          \item \al{甲暴力绑架乙后，乙反复向甲求情，甲释放了乙}。
          \item \al{甲收买 1 名儿童打算日后卖出，次日可怜他放走了}。
          \item \al{甲未能 xx，属未遂。主动 xx，属中止 \cuo}。
        \end{enumerate}
\end{enumerate}
